\subsection{Cómputo de Congruencias por Dígitos (Regla de Horner)}

\graybold{Preguntas guía:}

\begin{itemize}
\item ¿Necesito el residuo (mod m) de un número representado por dígitos, demasiado grande para operar como un entero normal en otros lenguajes?
\item ¿Puedo procesar los dígitos de izquierda a derecha manteniendo solo un residuo acumulado?
\item ¿La base y el módulo no tienen una relación especial (a diferencia de la regla de divisibilidad alternada, que aprovecha base $\equiv -1 \bmod m$)?
\end{itemize}

\graybold{¿Para qué sirve?} Calcular el residuo módulo $m$ de un número extremadamente grande (dado por sus dígitos en alguna base) sin construir el número completo, usando la regla de Horner.

\graybold{Complejidad:} \bigO{L}, con L = cantidad de dígitos.

\graybold{Fundamento teórico:} Si $N = d_0 d_1 \ldots d_{L-1}$ en base $b$, $N \bmod m$ se calcula incrementalmente: $r_0 = d_0 \bmod m$; $r_i = (r_{i-1} \cdot b + d_i) \bmod m$. Al terminar, $r_{L-1} = N \bmod m$. Esto generaliza a cualquier base y módulo, a diferencia de trucos específicos como la suma alternada (que solo aplica cuando base $\equiv -1 \bmod m$).

\graybold{Ejercicio aplicado}

Dado un número H expresado en binario (hasta 500 dígitos), calcular cuántos huevos sobran al empacarlos en cajas de 12 (es decir, $H \bmod 12$).

\graybold{I/O:} $|B| \le 500$ por línea, múltiples casos hasta el centinela '*' (patrón 3) → readline por línea es más que suficiente.

\graybold{Código solución}

\begin{lstlisting}[language=Python]
import sys

def solve():
    out = []
    for raw in sys.stdin:
        line = raw.strip()
        if line == '*':
            break
        rem = 0
        for ch in line:           # regla de Horner digito a digito
            rem = (rem * 2 + (1 if ch == '1' else 0)) % 12
        out.append(str(rem))
    sys.stdout.write("\n".join(out) + "\n")

solve()

# Nota: en Python, int(B, 2) % 12 tambien funcionaria directo (enteros de
# precision arbitraria), pero la tecnica de Horner generaliza a numeros
# aun mas grandes o a lenguajes sin enteros arbitrarios.
\end{lstlisting}
